\documentclass[10pt,a4paper]{article} 

\usepackage[left=1.5cm,right=1.5cm,top=1cm]{geometry}

\usepackage{enumitem}
\usepackage{array} 
\usepackage{tikz}
\usepackage{hyperref}
\usepackage{amssymb, amsmath}
\usepackage{fontawesome5}
\usepackage[spanish]{babel}

\newcommand{\CatLined}[1]{ \vspace{2em}\noindent\tikz{\draw (0,0)--(\linewidth, 0) node[pos=0.25,fill=white]{\bfseries{#1}}}}
\newcommand{\CatBold}[1]{ \vspace{2em} {\hfill {\bfseries{#1}} \hfill } }

\setlist[itemize]{label=\raisebox{0.25ex}{\tiny$\bullet$},topsep=1ex,itemsep=0pt}

\begin{document}
\fontfamily{cmss}\selectfont

\begin{center}
    {\Large\bfseries Camila Paz Guajardo Vásquez}\\
    \vspace*{0.25 cm}
    %{\href{https://github.com/cami-gv}{\raisebox{-0.05\height}\faGithub\ cami-gv} \ $|$ \ }
    {\href{https://cami-gv.github.io}{\faGlobe \ Camila Guajardo Vásquez} \ $|$ \ }
    {{\faEnvelope \ caguajardo@ug.uchile.cl} \ $|$ \ }
    {{\faMobile \ +569 5370 0715} }
\end{center}

Estudiante de Magíster en Ciencias Matemáticas interesada en topología geométrica, geometría hiperbólica y análisis geométrico.
Mi trabajo actual aborda la acción del grupo de automorfismos de una superficie de Riemann compacta sobre el conjunto de sus sístoles.

\CatLined{Educación}
\begin{itemize}
  \item
  {Magíster en Ciencias Matemáticas} \hfill \texttt{2024\,-\,presente}
  \\
  {\itshape Universidad de Chile}
  \item
  {Licenciatura en Matemáticas} \hfill \texttt{2020\,-\,2023}
  \\
  {\itshape Pontificia Universidad Católica de Chile}
\end{itemize}


\CatLined{Experiencia en investigación}
\begin{itemize}
    \item {Tesista en Proyecto Fondecyt 12340034} \hfill \texttt{2024}
    \newline \vspace*{0.25 cm}
    \begin{tabular}{rl}
        Tema:& Acciones de grupo en superficies de Riemann compactas  y sus efectos sobre \\
        &\quad el conjunto de sus geodésicas cerradas\\
        Investigadora responsable:& Prof. Anita Rojas
    \end{tabular}

    \item {Taller de Iniciación Científica} \hfill \texttt{2023}
    \newline\vspace*{0.25cm} \hspace*{0.5 cm}
    \begin{tabular}{rl}
        Tema:& \emph{Mapping Class Group} de Superficies \\
        Supervisor:& Prof. Mauricio Bustamante Londoño
    \end{tabular}

    \item {Investigaciones de Pregrado (Vicerrectoría de Investigación, PUC)} \hfill \texttt{2022-2023}
    \vspace*{-0.15 cm}
    \begin{itemize}
        \item La fórmula de Riemann-Hurwitz y sus aplicaciones, con Prof. Natalia García Fritz 
        \item Teoría de Morse y aplicaciones en topología, con Prof. Mauricio Bustamante Londoño
        \item Geometría de curvas planas en grados bajos, con Prof. Giancarlo Urzúa 
    \end{itemize}

\end{itemize}

\CatLined{Premios y reconocimientos}
\begin{itemize}
    \item {Beca de Magíster Nacional ANID} \hfill \texttt{2025}
    \item {Premio de Excelencia Académica de la generación 2023, Licenciatura en Matemáticas} \hfill \texttt{otorgado en 2024}
\end{itemize}

\CatLined{Charlas y seminarios selectos}
\begin{itemize}
    \item {Organizadora, Seminario de lectura en Geometría Riemanniana} \hfill \texttt{2025-2}
    \newline \vspace*{0.25 cm}
    \begin{tabular}{rl}
        Temas:& Variedades suaves, espacios tangentes, métricas Riemannianas, espacios modelo, \\
        &\quad conexiones, geodésicas, y curvatura\\
        Asesor:& Prof. Pedro Gaspar 
    \end{tabular}

    \item {Seminario No-degenerado, Universidad de Talca} \hfill \texttt{2025}
    \newline \vspace*{0.25 cm}\hspace*{0.25 cm} 
    Tema: Cotas para sístoles a partir de polígonos hiperbólicos

    \item {Seminario Relajao', Universidad de Chile} \hfill \texttt{2024-2025}
    \newline \vspace*{0.25 cm}\hspace*{0.25 cm} Temas: Superficies \emph{concretas} usando SageMath; Introducción a la teoría de Morse

    \item {Seminario de Teoría de Números, PUC} \hfill \text{2022-2023}
    \newline \vspace*{0.25cm} \hspace*{0.25cm} 
    Temas: Sobre fibras aditivas y superficies elípticas sobre \(\mathbb{P}^1\);
    Curvas elípticas y sus puntos racionales
    
    %\item {Seminario de Grupos Formales, PUC} \hfill \texttt{2022}
    %\newline \vspace*{0.25cm}\hspace{0.25cm} 
    %Temas: Ley de grupo formal universal; Grupo formal asociado a una curva elíptica

    %\item {Seminario de Combinatoria, PUC} \hfill \texttt{2022}
    %\newline \vspace*{0.25cm} \hspace*{0.25cm}
    %Tema: Sobre el paquete de Kähler de un anillo

\end{itemize}

\newpage
\CatLined{Cursos optativos selectos}
\begin{itemize}
    \item {POST664\,--\,Análisis II con Prof. Nikola Kamburov} \hfill \texttt{6.4/7.0}
    \newline \vspace*{0.25cm}
    \begin{tabular}{rl}
    Contenidos:& Espacios de Banach, teoremas de Baire y Hahn--Banach, espacios de Hilbert \\
    & y teorema de Lax-Milgram, teoría espectral de operadores compactos autoadjuntos, \\
    & espacios de Sobolev, teoría de distribuciones
    \end{tabular}

    \item {POST5900\,--\,Introducción a las superficies de Riemann y curvas algebraicas con Prof. Anita Rojas} \hfill \texttt{7.0/7.0}
    \newline \vspace*{0.25cm}
    \begin{tabular}{rl}
    Contenidos:& Superficies de Riemann, curvas algebraicas en \(\mathbb{P}^2(\mathbb{C})\), funciones meromorfas, \\
    & aplicaciones holomorfas entre superficies de Riemann, fórmula de Riemann-Hurwitz 
    \end{tabular}

    \item {CSDMATE157\,--\,Toros complejos con Prof. Robert Auffarth} \hfill \texttt{6.1/7.0}
    \newline \vspace*{0.25cm}
    \begin{tabular}{rl}
    Contenidos:& Homomorfismos de toros, grupo de Néron--Severi, toro dual, extensiones, \\
    & complementos de subtoros, toros simples e indescomponibles
    \end{tabular}

    \item {CS07MM0019\,--\,Topología y Geometría de Variedades con Prof. Mauricio Bustamante} \hfill \texttt{6.7/7.0}
    \newline \vspace*{0.25cm}
    \begin{tabular}{rl}
    Contenidos:& Variedades suaves, espacios tangentes, diferencial de aplicaciones suaves,\\
    & inmersiones y submersiones, teorema de Sard, transversalidad
    \end{tabular}

    \item {MAT2555\,--\,Análisis Funcional con Prof. Nikola Kamburov} \hfill \texttt{6.3/7.0}
    \newline \vspace*{0.25cm}
    \begin{tabular}{rl}
    Contenidos:& Espacios de Banach, teoremas de Baire y de Hahn--Banach, espacios de Hilbert, \\
    & series de Fourier, teoría espectral de operadores compactos autoadjuntos
    \end{tabular}

    \item {MAT2850\,--\,Topolog\'ia Algebraica con Prof. Siddarth Mathur} \hfill \texttt{6.3/7.0}
    \newline \vspace*{0.25cm}
    \begin{tabular}{rl}
    Contenidos:& Espacios de cubrimiento, grupo fundamental; homología singular, relativa y celular
    \end{tabular}

    %\item {MAT2335\,--\,Introducción a la Geometría Algebraica con Prof. Natalia García} \hfill \texttt{6.0/7.0}
    %\newline \vspace*{0.25cm}
    %\begin{tabular}{rl}
    %Contenidos:& Curvas algebraicas, variedades afines y proyectivas, morfismos birracionales y dominantes
    %\end{tabular}

    %\item {MAT2225\,--\,Teoría de Números con Prof. Ricardo Menares} \hfill \texttt{6.1/7.0}
    %\newline \vspace*{0.25cm}
    %\begin{tabular}{rl}
    %Contenidos:& Teorema fundamental de la aritmética, ecuaciones diofantinas, teorema chino de los restos,\\
    %& ley de reciprocidad cuadrática
    %\end{tabular}

    \item {MAT2305\,--\,Geometría Diferencial, con Prof. Mauricio Bustamante} \hfill \texttt{6.4/7.0}
    \newline \vspace*{0.25cm}
    \begin{tabular}{rl}
    Contenidos:& Curvas en \(\mathbb{R}^3\), teoría local de superficies en \(\mathbb{R}^3\) (formas fundamentales, curvaturas y geodésicas), \\
    & teorema de Gauss-Bonnet
    \end{tabular}

\end{itemize}


\CatLined{Experiencia en docencia}
\begin{itemize}
    \item {Ayudante de cátedra y corrección. Facultad de Matemáticas, Pontificia Universidad Católica de Chile}
    \vspace*{-0.15cm}
    \begin{itemize}
        \item MPG3403\,--\,Topología y Geometría de Variedades, Prof. Pedro Gaspar \hfill \texttt{2025}
        \item MAT2704/MAT270I\,--\,Variable Compleja, Prof. Martin Chuaqui \hfill \texttt{2024}
        \item MAT1389\,--\,Geometría I, Prof. María José Moreno y Prof. María Fernanda Espinal \hfill \texttt{2023}
        \item MAT2218\,--\,Álgebra abstracta II, Prof. Ricardo Menares \hfill \texttt{2023}
        \item MAT1204\,--\,Introducción al Álgebra, Prof. Amal Taarabt \hfill \texttt{2022}
    \end{itemize}

    \item {Ayudante de cátedra y corrección. Departamento de Matemáticas, Universidad de Chile}
    \vspace*{-0.15cm}
    \begin{itemize}
        \item CVER004\,--\,Álgebra y Geometría II, Prof. Anita Rojas \hfill \texttt{2025}
        \item MCLM110/MC110\,--\,Álgebra y Geometría I, Prof. Nelda Jaque Tamblay \hfill \texttt{2024}
    \end{itemize}

    %\item {Apoyo al aprendizaje. Universidad de Chile}
    %\vspace*{-0.15cm}
    %\begin{itemize}
    %    \item Sala de ayuda de la Escuela de Ciencias Ambientales y Biotecnología \hfill \texttt{2025}
    %\end{itemize}

    %\item {Apoyo al aprendizaje. Pontificia Universidad Católica de Chile}
    %\vspace*{-0.15cm}
    %\begin{itemize}
    %    \item MAT0004\,--\,Taller de Matemáticas \hfill \texttt{2022}
    %    \item Sala de ayuda en Matemáticas \hfill \texttt{2021}
    %    \item MAT0006\,--\,Taller de Matemáticas \hfill \texttt{2021}
    %\end{itemize}

\end{itemize}


%\CatLined{Participación en proyectos}
%\begin{itemize}
%    \item {Revisora y editora de repositorio de experimentos matemáticos. Fondo de Incentivo a la Investigación de la Docencia en Pregrado. Universidad de Chile} \hfill \texttt{2024}
%    \begin{tabular}{rl}
%        Título del proyecto:& Diseño de experimentos matemáticos como parte de la reflexión de la práctica docente \\
%        & para contribuir a una mejora en la trayectoria formativa de estudiantes en cursos\\
%        & de matemáticas\\
%        Investigadoras:& Anita Rojas, Leslie Jiménez, Carolina Canales, Natalia Henríquez y Cristóbal Rivas
%   \end{tabular}
%\end{itemize}


\CatLined{Divulgación y servicio}
\begin{itemize}
    \item {Monitora, Festival de Matemáticas en Puerto de Ideas} \hfill \texttt{Antofagasta, 2023}
    \item {Entrenadora, Olimpiadas Internacionales de Matemáticas} \hfill \texttt{2021\,-\,presente}
    \item {Correctora, Olimpiada Nacional de Matemáticas} \hfill \texttt{2021\,-\,presente}
    \item {Miembro, Equipo Académico CMat} \hfill \texttt{2020\,-\,2022}
    \item {Coordinadora de Problema, I Olimpiada Panamericana Femenil de Matemáticas} \hfill \texttt{2021}
    
\end{itemize}

\end{document}

% ¿Poner las participaciones olímpicas?
